% ============================================================================
% THE GPU BESTFILL PIPELINE
% ============================================================================
\subsection{Problem Statement}
\label{sec:bestfill_problem}

Given a context chord $\mathbf{q} \in \{0,1\}^D$ and a
dictionary (PrimeField) of $N$ stored IcosahedralManifolds
$\{\mathcal{D}_i\}_{i=1}^{N}$, the BestFill operation returns:

\begin{equation}
  i^* = \operatorname*{arg\,max}_{1 \leq i \leq N}\;
        \text{Score}(\mathbf{q}, \mathcal{D}_i)
  \label{eq:bestfill_argmax}
\end{equation}

Subject to the constraint that $\mathbf{q}$ and $\mathcal{D}_i$ pass all
pre-filter checks.  This is the system's sole inference primitive---all
generation, retrieval, and reasoning reduce to repeated BestFill calls.

\subsection{The 5-Pass Retrieval Pipeline}
\label{sec:five_pass}

The GPU executes a hierarchical elimination pipeline designed to exploit
SIMT warp architecture.  Each pass is strictly monotone:
$N_{\text{pass}_{k+1}} \leq N_{\text{pass}_k}$.

\paragraph{Pass 1: Winding Filter (Integer, $O(1)$ per candidate).}
Extract the 4-bit winding counter $w_i$ from header:
\begin{equation}
  w_i = \bigl(\text{Header}_i \gg 5\bigr) \mathbin{\&} 0\text{xF}
  \label{eq:winding}
\end{equation}
Reject candidates where $w_i \neq w_q$.  This eliminates manifolds from
incompatible syntactic depths.

\paragraph{Pass 2: Group State Filter (Integer, $O(1)$).}
Extract the 1-bit state flag $\sigma_i = \text{Header}_i \gg 15$.
Reject candidates where $\sigma_i \neq \sigma_q$ (Cubic $\neq$
Icosahedral).

\paragraph{Pass 3: Coarse Popcount (Bitwise SIMD, $O(D/w)$).}
Compute the condensed macro-shadow chord $\hat{\mathbf{d}}_i$
(ChordOR over all 257 blocks of \texttt{Cubes[0]}) and evaluate:

\begin{equation}
  \hat{s}_i = \texttt{popcount}(\mathbf{q} \mathbin{\&} \hat{\mathbf{d}}_i)
  \label{eq:coarse_score}
\end{equation}

Candidates with $\hat{s}_i < \tau_{\text{coarse}}$ are discarded before
the expensive full-structure comparison.

\paragraph{Pass 4: Dense Popcount (Memory Bandwidth, $O(135 \cdot D/w)$).}
For surviving candidates, compute the full FillScore over all active blocks:

\begin{equation}
  \text{Score}_i =
    \frac{\sum_{c=0}^{4}\sum_{j=0}^{26}
          \texttt{popcount}\bigl(\mathbf{q} \mathbin{\&}
                                \mathcal{D}_i.\text{Cubes}[c][j]\bigr)}
         {\sum_{c,j}
          \texttt{popcount}\bigl(\mathbf{q} \mathbin{\&}
                                \mathcal{D}_i.\text{Cubes}[c][j]\bigr)
          \;+\;
          \sum_{c,j}
          \texttt{popcount}\bigl(\mathcal{D}_i.\text{Cubes}[c][j]
                                \mathbin{\&} \lnot\,\mathbf{q}\bigr)
          \;+\; 1}
  \label{eq:dense_score}
\end{equation}

This is the multi-block generalization of FillScore
(\Cref{eq:fill_score}), summing resonance and noise across all 135
chord fields.  The $+1$ in the denominator prevents division by zero.

\paragraph{Pass 5: Geodesic LUT Disambiguation (L1 Cache, $O(1)$).}
For tied scores, the 6-bit \texttt{RotState} indices $r_i, r_j$ of
competing candidates are used to fetch exact geometric distances from the
precomputed $60 \times 60$ Unified Geodesic Matrix:

\begin{equation}
  d(r_i, r_j) = \arccos\bigl(|q_i \cdot q_j|\bigr), \qquad
  d \in [0, \pi]
  \label{eq:geodesic_lut}
\end{equation}

The candidate with shorter geodesic distance to the query state wins.
The entire matrix occupies 3,600 bytes and fits in L1 cache, making
lookups free relative to memory bandwidth.

\subsection{Packed Score Representation}
\label{sec:packed_score}

To enable lock-free parallel reduction across GPU threads, the score
and index are packed into a single atomic \texttt{uint64}:

\begin{equation}
  \text{packed} = \bigl(\lfloor \text{Score} \times 2^{32}\rfloor \ll 32\bigr)
                  \;\big|\; \text{Index}
  \label{eq:packed}
\end{equation}

The winner is determined by a single \texttt{AtomicMax} operation per
thread group.  Because the score occupies the high 32 bits, lexicographic
integer comparison correctly preserves score ordering.

\subsection{Multi-Backend Dispatch}
\label{sec:multi_backend}

The kernel dispatch layer supports Metal (Apple Silicon), CUDA (NVIDIA),
and CPU (fallback) backends with automatic fallback ordering.  For large
dictionaries ($N > N_{\text{shard}}$), heterogeneous local sharding
distributes work:

\begin{equation}
  i^* = \operatorname*{arg\,max}_{b \in \text{Backends}}
        \operatorname*{arg\,max}_{i \in \text{Shard}(b)}
        \text{Score}(\mathbf{q}, \mathcal{D}_i)
  \label{eq:hetero_shard}
\end{equation}

Each backend atomically publishes its local winner via
\texttt{CompareAndSwap} on the global packed score.

\subsection{Expected Reality Modulation}
\label{sec:expected_reality}

BestFill accepts an optional \emph{ExpectedReality} manifold
$\mathbf{e} \in \{0,1\}^D$---a structural template representing the
anticipated shape of the continuation.  When provided, the vacuum is
computed against $\mathbf{e}$ rather than $\mathbf{q}$:

\begin{equation}
  \mathbf{h} = \mathbf{e} \mathbin{\&} \lnot\,\mathbf{q}
  \label{eq:expected_reality}
\end{equation}

and FillScore is evaluated as $\text{FillScore}(\mathbf{h}, \mathcal{D}_i)$,
enabling the system to search for spans that fill a specific structural
vacuum rather than merely maximizing context overlap.
